\section{Introduction}

While establishing classification theory, Shelah introduced the eq construction of a first-order theory $T$ and studied what he called imaginary elements \cite[Ch.~III \S6]{shelahbook}. This was done so that every set definable with parameters has a minimal set of defining parameters, called its canonical parameter. The idea is simple: an imaginary element is an equivalence class of a definable equivalence relation; the eq construction expands a theory $T$ to $T^{eq}$ by freely adding in quotients of definable equivalence relations as new sorts. In \cite{poizat_imaginary_galois}, Poizat defined that $T$ eliminates imaginaries if such quotients can be naturally identified with a definable set in $T$. He went on to show that the theory of algebraically closed fields and the theory of differentially closed fields with characteristic zero both eliminate imaginaries, thereby establishing a model-theoretic generalisation of classical Galois theory and differential Galois theory. More recently, Hrushovski \cite{hrushovski_groupoid} studied generalised imaginaries, replacing definable equivalence relations with definable groupoids, which are the vertical categorification of equivalence relations. He showed that generalised imaginaries are related to internal covers of $T$ and type amalgamation. This is subsequently generalised in \cite{GoodrickJohn2013Afab,GoodrickJohn2014Tpap,GoodrickJohn2010Gca3,kamensky_cat_app_int,kamensky_higher_int_cov,moosa_haykazyan,wang_simp_grpd}, providing a fecund field for interactions between model theory and category theory, and indeed higher category theory.

\smallskip

On the other hand, category theorists were also studying equivalence relations internal to a category and the existence of their quotients, as well as the free completion of a category under these quotients. A category with all such quotients is called exact\footnote{There are two unrelated notions of exactness in category theory, called Barr-exactness and Quillen-exactness. The exactness considered in this paper is Barr-exactness.}; the free construction is called the ex/reg completion; we write $\C_{ex/reg}$ for the ex/reg completion of a regular category $\C$. The subscript `ex' stands for exact and `reg' stands for regular. The ex/reg completion freely turns a regular category into an exact category. This is to distinguish it from the other exact completion---the ex/lex completion---which freely turns a lex category---a category with finite limits---into an exact category. The first references come from Barr \cite{Barrbook} and Lawvere \cite{Lawvere_Perugia}. 
It is evident that Lawvere had already known about what was eventually named ex/reg completion. Indeed, in \cite[pp.~84--85]{Lawvere_Perugia}) he sketched out the construction, but it would be Succi Cruciani \cite{cruciani}, and Carboni and Vitale \cite{CarboniA.1998Raec}, who proved that the construction described by Lawvere is indeed the free construction and has the expected universal property. Lack \cite{LackStephen1999Anot} provides another description for this construction in terms of sheaves, which was later generalised by Shulman \cite{ShulmanMichael2012Ecas}. There is a related construction---the pretopos completion of a coherent category---which often coincides with the ex/reg completion; for example, see \cite[Proposition 5.1]{EmmeneggerJacopo2020Edac}. As a result, there is potential for confusion for which construction is relevant in which setting. Shulman's work \cite{ShulmanMichael2012Ecas} provides a framework of seeing the ex/reg completion as the finitary version of the pretopos completion.

\smallskip

The connection between the model-theoretic approach and the category-theoretic approach is attributed to Makkai, who, according to Harnik's expository paper \cite{harnik}, knew already in 1980 that `$T^{eq}$ is nothing but the pretopos completion'. Indeed, the subject of categorical model theory was established in Makkai's joint work with Reyes in \cite{alma992977584040101631}. Briefly, given a complete first-order theory $T$, one can define the syntactic category of $T$ as the category of definable sets and definable functions, denoted by $\defT(T)$. Since our logic is first-order, it is in particular regular, which is the fragment of first-order logic consisting of truth $\top$, finite conjunction $\wedge$, equality $=$, and existential quantification $\exists$. Hence, $\defT(T)$ is a regular category. In fact, the inclusion of finite disjunction $\vee$ in the logic makes $\defT(T)$ coherent, and the inclusion of negation $\neg$ makes $\defT(T)$ Boolean. The syntactic category encodes all the information of $T$ in a language-agnostic way: the category does not `see' whether a formula is quantifier-free or not. Nonetheless, much model theory can be done by studying $\defT(T)$; for example, the category $\textup{Coh}(\defT(T),\textbf{Set})$ of coherent functors into the category of sets is equivalent to the category of models and elementary embeddings \cite[Ch. 8 \S1]{alma992977584040101631}. 

\smallskip 

The work of Makkai and Reyes includes the celebrated result known as conceptual completeness \cite[Ch. 7]{alma992977584040101631}. This says that if an interpretation of a theory $T_1$ in another theory $T_2$ induces an equivalence between the two categories of models, then this interpretation is in fact a bi-interpretation up to the eq construction, or equivalently, up to the pretopos completion. This is perhaps another reason why the eq construction has become synonymous with the pretopos completion, as seen in \cite{harnik}. In \S\ref{preliminaries}, we will review this and will find that, it is exactness, rather than the property of being a pretopos, that is relevant to imaginaries. Indeed, we see in Theorem~\ref{degenerate theorem} that certain non-degenerate assumptions are needed for the pretopos completion of $\defT(T)$ to coincide with $\defT(T^{eq})$. On the contrary, as we shall see in Theorem \ref{Teq is exact}, no such non-degenerate assumptions are needed to prove 
\[\defT(T^{eq})\simeq\defT(T)_{ex/reg}.\]

In \S\ref{pro}, we study the pro-completion $\Pro(\defT(T))$ of $\defT(T)$ with a view of generalising in \S\ref{hyperimaginaries} the connection between imaginaries and exactness to a connection between the projective limits of imaginaries, called hyperimaginaries, and the exactness of $\Pro(\defT(T))$. While instances of projective limits, or equivalently cofiltered limits, have been studied for over a century---Hensel was studying the $p$-adic numbers as early as 1897---the completion $\Pro (\C)$ of a category $\C$ under projective limits was first introduced by Grothendieck \cite{Grothendieck1958-1960}. Similar to the ex/reg completion, the pro-completion freely adds projective limits into a category. This completion can be described as a certain subcategory of the category of copresheaves on $\C$. When one dualises this construction, one obtains the ind-completion, which freely adds in direct limits, or equivalently filtered colimits \cite[Expos\'{e} I \S8.10]{a.TheorieToposCohomologie1972}. 

\smallskip

Barr \cite{BARR1986113} studied $\Pro (\C)$ in the case when $\C$ is regular. It turns out that many nice properties of $\C$ transfer to $\Pro(\C)$, such as regularity and coherence. Some properties do not transfer, for example, being Boolean. Model-theoretically, this is saying that the complement of a type-definable set is a $\bigvee$-type definable set, that is, definable with infinitary disjunction, rather than a type-definable set. Another property that does not transfer is exactness. 

\begin{example}\label{finset}
    Let $\C$ be the category \textbf{FinSet} of finite sets and all functions, which is a pretopos and in particular exact. Its pro-completion is the category \textbf{Stone} of Stone spaces and continuous functions (see \cite[Ch.~VI \S 2.3]{alma9912295854401631}), which is not exact. Indeed, its ex/reg completion is \textbf{CHaus}, the category of compact Hausdorff spaces and continuous functions; see \cite[Theorem 8.1]{MarraVincenzo2020Acot}. This fits into the $\defT(T)$ framework: consider $T$ to be the theory of the structure with sorts $S_n$ for each $n\in\omega$, where each sort $S_n$ is interpreted as a set containing exactly $n$ elements, with every element named by a constant. Then, $\defT(T) \simeq \textbf{FinSet}$.
\end{example}

Of course, there are examples where $\Pro(\C)$ is exact. When $\C$ is the category \textbf{FinGrp} of finite groups, the pro-completion is the category \textbf{ProfGrp} of profinite groups and continuous group homomorphisms. Now, equivalence relations of profinite groups correspond to closed subgroups. Since the quotient of a profinite group by a closed subgroup is again a profinite group \cite[Proposition~2.2.1~(a)]{profinite_groups}, \textbf{ProfGrp} is exact.

\smallskip

It was Kamensky \cite{KAMENSKY2007180}, who first studied $\Pro(\defT(T))$ in the context of model theory, relating it to the types of $T$. By types, we mean $\ast$-types without parameters, which are sets of formulas in an infinite tuple of variables; likewise, given a model $\M$, a type-definable set is a subset of $\M^I$, whose elements are the realisations of a type and where $I$ is an infinite indexing set. Then, the objects of $\Pro(\defT(T))$ can be thought of as the types of $T$. The morphisms correspond to types, whose interpretation in every model is the graph of a function. Equivalently, one can fix a monster model $\U$ and ask for the interpretation in $\U$ to be the graph of a function. In \S\ref{subsect_saturation}, we show a new characterisation of saturation in terms of the preservation of certain regular epimorphisms in $\Pro(\defT(T))$; namely, a model $\M$ is $\kappa$-saturated if and only if for any regular epimorphism in $\Pro(\defT(T))$, whose codomain has variable length strictly less than $\kappa$, the interpretation in $\M$ is a surjection. Now, hyperimaginaries are roughly speaking quotients of types. If we wish to study them categorically, we would need a category of types. Hence, $\Pro(\defT(T))$ is exactly the right setting to study hyperimaginaries.

\smallskip

Although stability is arguably the poster child of classification theory, a slightly weaker property---simplicity---introduced by Shelah \cite{SHELAH1980177}, also received attention for its nice behaviours, culminating in the celebrated Kim-Pillay Theorem \cite{KIM1997149}. One thing that model-theorists noticed was that, while working in stable theories usually requires one to deal with imaginaries, working in simple theories requires one to deal with hyperimaginaries, which are heuristically `type-definable imaginaries'. More precisely, a hyperimaginary is an equivalence class of a type-definable equivalence relation. We say a theory eliminates hyperimaginaries if every hyperimaginary is interdefinable with a sequence of imaginaries; this is equivalent to saying that every type-definable equivalence relation is the intersection of some definable equivalence relations. The first mention of this is in \cite{Pillay_Poizat_1987}, where Pillay and Poizat proved that stable theories eliminate hyperimaginaries. In \cite{BuechlerSteven2001St}, this was generalised to supersimple theories, yet whether simple theories eliminate hyperimaginaries remains an open question. 

\smallskip

So far, the model theory and the category theory are well-known, though perhaps only to model theorists and categorical theorists respectively. In order to make this paper accessible to both communities, we include basic definitions, as well as proofs for elementary facts. We hope this will be helpful in further bridging the two areas.

\smallskip

We come to the main original contributions of the paper in \S\S\ref{hyperimaginaries}, \ref{heq}. We are interested in whether $\Pro(\defT(T))$ is exact. Now, we shall see in Proposition~\ref{if prodef exact then def exact} that, if the quotient of a definable equivalence relation is type-definable, then the quotient itself is also definable. In other words, if we want $\Pro(\defT(T))$ to be exact, then we should require $\defT(T)$ to also be exact. Hence, we might as well assume $T$ eliminates imaginaries, or work over $\defT(T^{eq})$, equivalently $\defT(T)_{ex/reg}$. Assuming this, elimination of hyperimaginaries allows us to decompose a type-definable equivalence relation as a cofiltered limit of definable equivalence relations, each of which has a definable quotient. This allows us to obtain in Theorem~\ref{eliminate hyper iff pro exact} the following characterisation:
\[T \textup{ eliminates hyperimaginaries } \Leftrightarrow \ \Pro(\defT(T)_{ex/reg}) \textup{ is exact.}\]

Similar to the eq construction for imaginaries, there is an heq construction for hyperimaginaries. We study this in \S\ref{heq}. The construction comes from \cite{Hart2000-HARCAC-24}, but unlike the eq construction, hyperimaginaries cannot be na\"{i}vely added to a structure as real elements. Indeed, \cite[Example 3.1.6]{alma992983361725701631} shows that the quotient of a type-definable equivalence relation can have bounded infinite cardinality. If we simply add the quotient as a new sort, by compactness there are elementarily equivalent structures, whose interpretations of that sort have unbounded cardinalities. Moreover, one cannot add the quotient map as a new function symbol, since the sort of the domain might be of infinite length. Nevertheless, \cite{Hart2000-HARCAC-24} constructs $\U^{heq}$ from a monster model $\U$ as an expansion with quotients of type-definable equivalence relations as new sorts. Instead of considering $\U^{heq}$ as a first-order structure in a certain language, one considers how automorphisms of $\U$ extends to automorphisms of $\mathcal{U}^{heq}$ and, using this, defines the various closure operators like dcl, acl and bbd.

\smallskip

Shortly afterwards, Ben-Yaacov started developing positive model theory and noted in \cite[Example 2.16]{Ben-Yaacov2003-BENPMT} that $T^{heq}$ can be given a syntax. This was further developed by Dobrowolski and Kamsma \cite{DobrowolskiJan2022Kipl} to study NSOP$_1$ in positive logic. Omitting negation, as in positivie logic, turns out to be the correct choice for studying hyperimaginaries, since type-definable sets are not closed under complement. In \S\ref{heq}, we will argue that the correct logic for studying hyperimaginaries is a logic with infinite conjunctions, infinite existential quantification, finite disjunction, and no negation. This is the internal logic of $\infty$-coherent categories, such as $\Pro(\defT(T))$. Using this, we will define an infinitary syntactic category associated to $\U^{heq}$, denoted by $\tpdefT(\U^{heq})$, and prove that 
\[\tpdefT(\U^{heq}) \simeq (\Pro(\defT(T)))_{ex/reg}. \] 

We summarise our results.

\begin{mainthm}
    
    Let $T$ be a consistent complete first-order theory and let $\U$ be a monster model.
    \begin{enumerate}
    \item $T$ eliminates imaginaries if and only if $\defT(T)$ is exact.

    \item $T$ eliminates hyperimaginaries if and only if $\Pro(\defT(T)_{ex/reg})$ is exact.
        
    \item $T$ eliminates imaginaries and hyperimaginaries if and only if $\Pro(\defT(T))$ is exact.

    \item We have the following equivalences of categories: 
        \begin{align*}
            \defT(T^{eq}) &\simeq \defT(T)_{ex/reg}, \\
            \tpdefT(\U^{heq}) &\simeq (\Pro(\defT(T)))_{ex/reg}.
        \end{align*}
    \end{enumerate}


\end{mainthm}

\smallskip

\noindent\textbf{Acknowledgements.} This work is part of my PhD research, done under the supervision of Omar Le\'on S\'anchez and the co-supervision of Nicola Gambino. I would like to thank both for their guidance. I would also like to thank Mark Kamsma, who posed the question that this paper is answering, as well as Giacomo Tendas for helpful discussions. 
